% Generated by tables/build_tables.py; do not edit.
\begin{tabular*}{\textwidth}{@{\extracolsep{\fill}}llcrcrrr@{}}
\toprule
 &  &  &  & First test & Mean & Volatility & Autocorr. \\
Series & Source & Span & Days & month & (\% p.a.) & (\% p.a.) & (lag 1) \\
\midrule
US market & Ken French & 1926--2026 & 26,338 & 1930-10 & 7.8 & 17.1 & 0.046 \\
Developed markets & Ken French & 1990--2026 & 9,458 & 1995-05 & 6.4 & 14.4 & 0.121 \\
North America & Ken French & 1990--2026 & 9,458 & 1995-05 & 9.0 & 17.5 & $-$0.056 \\
Europe & Ken French & 1990--2026 & 9,458 & 1995-05 & 6.2 & 17.7 & 0.031 \\
Japan & Ken French & 1990--2026 & 9,458 & 1995-05 & 2.9 & 21.3 & $-$0.035 \\
Asia Pacific ex Japan & Ken French & 1990--2026 & 9,458 & 1995-05 & 6.8 & 16.3 & 0.124 \\
S\&P 500 (price) & FRED & 2016--2026 & 2,512 & 2021-12 & 14.1 & 18.1 & $-$0.147 \\
\bottomrule
\end{tabular*}
\par\vspace{4pt}
\begin{minipage}{\textwidth}\footnotesize\setstretch{1.0}\emph{Notes.} Daily total returns in US dollars (the S\&P 500 series is a price index without dividends, and its mean is a raw return). Mean is the annualised average return above the one-month T-bill. Autocorr.\ is the first-order autocorrelation of daily returns: positive values in regional indices mostly reflect stale prices across time zones, a known source of spurious predictability. The first test month leaves room for the longest training window (1,008 days), the gap and a 250-day burn-in.
\end{minipage}
